\documentclass[a4paper,ukrainian,oneside]{extreport}

% - Language & encoding -
\usepackage{cmap}
\usepackage[T2A]{fontenc}
\usepackage[utf8]{inputenc}
\usepackage[russian,ukrainian]{babel}

% - Math -
\usepackage{amssymb,amsfonts,amsmath,amsthm,amscd}
\usepackage{upgreek}
\usepackage{textcomp}
\usepackage{mathtext}
\usepackage{hanging}
\usepackage{etoolbox}
\usepackage{flafter}
\usepackage[shortlabels]{enumitem}

% - Page geometry (from bachelor_thesis.cls) -
\usepackage[includehead, headsep=0pt, footskip=0pt,
            top=2cm, bottom=2cm, left=3cm, right=1.5cm]{geometry}

% - Spacing & formatting -
\usepackage{indentfirst}
\usepackage{setspace}
\usepackage{fancyhdr}
\usepackage{tabularx}
\usepackage{csquotes}
\usepackage[titles]{tocloft}
\usepackage[hidelinks]{hyperref}

% - Chapter/section headings (bachelor_thesis.cls) -
\usepackage[rm,small,raggedright]{titlesec}

% - Figures & captions -
\usepackage{graphicx}
\usepackage[singlelinecheck=off]{caption}
\DeclareCaptionLabelFormat{tablecap}{\hspace*{1.25cm}Таблиця #2}
\DeclareCaptionLabelFormat{figurecap}{Рисунок #2}
\captionsetup[table]{margin=0pt, singlelinecheck=false, labelfont=bf, textfont=normalfont, justification=justified, labelsep=endash, labelformat=tablecap, font = {Large, stretch = 1.5}}
\captionsetup[figure]{justification=centering, labelsep=endash, labelfont=bf, textfont=normalfont, labelformat=figurecap, font = {Large, stretch = 1.5}}
\setlength{\textfloatsep}{21pt}
\setlength{\intextsep}{21pt}

% - Suppress word hyphenation (from cls) -
\tolerance=1
\emergencystretch=\maxdimen
\hyphenpenalty=10000
\hbadness=10000

\sloppy

% - Font size command -
\newcommand{\setfontsize}[1]{\fontsize{#1pt}{#1pt}\selectfont}

% - Paragraph indent -
\setlength{\parindent}{1.25cm}
\AddEnumerateCounter{\asbuk}{\@asbuk}{\cyrm}
\setlist[itemize,1]{label=--, partopsep=0pt, parsep=0pt, topsep=0pt, itemsep=0pt, itemindent=1.7cm, leftmargin=0pt}
\setlist[itemize,2]{label=--, partopsep=0pt, parsep=0pt, topsep=0pt, itemsep=0pt, itemindent=2.25cm, leftmargin=0pt}
\setlist[enumerate,1]{label=\arabic*), partopsep=0pt, parsep=0pt, topsep=0pt, itemsep=0pt, itemindent=1.9cm, leftmargin=0pt}
\setlist[enumerate,2]{label=\asbuk*), partopsep=0pt, parsep=0pt, topsep=0pt, itemsep=0pt, itemindent=2.6cm, leftmargin=0pt}

% - Page style: page number top-right -
\fancypagestyle{plain}{
    \fancyhf{}
    \rhead{\thepage}
    \renewcommand{\headrulewidth}{0pt}
    \renewcommand{\footrulewidth}{0pt}
}

% - Chapter format (bachelor_thesis.cls) -
\titleformat{\chapter}[block]
  {\centering\Large\bfseries}
  {РОЗДІЛ~\thechapter.}
  {0.5em}
  {\MakeUppercase}
\titlespacing*{\chapter}{0pt}{-28pt}{48pt}

% - Section format (bachelor_thesis.cls) -
\titleformat{\section}[block]
  {\Large\bfseries}
  {\hspace*{1cm}\thesection}
  {0.5em}
  {}
\titlespacing{\section}{0.5em}{24pt}{24pt}

% - Subsection format (bachelor_thesis.cls) -
\titleformat{\subsection}[block]
  {\Large\bfseries}
  {\thesubsection}
  {0.5em}
  {}
\titlespacing{\subsection}{\parindent}{24pt}{24pt}

\preto\section{%
  \ifnum\value{section}=0\addvspace{-24pt}\fi%
}
\preto\subsection{%
  \ifnum\value{subsection}=0\addvspace{-24pt}\fi%
}

% - Numbering per chapter -
\renewcommand{\theequation}{\thechapter.\arabic{equation}}
\renewcommand{\thefigure}{\thechapter.\arabic{figure}}

% - Зміст customisation (from bachelor_thesis.cls) -
\makeatletter
\renewcommand{\tableofcontents}{
    \clearpage
    \thispagestyle{empty}
    \begin{center}
        \bfseries\MakeUppercase{Зміст}
    \end{center}
    \addvspace{42pt}
    \@starttoc{toc}
}
\makeatother

% dot leaders
\renewcommand{\cftchapdotsep}{1}
\renewcommand{\cftsecdotsep}{1}
\renewcommand{\cftsubsecdotsep}{1}
\cftsetpnumwidth{1em}
\setlength\cftbeforechapskip{0pt}

% chapter entries
% chapter entries
\renewcommand{\cftchapaftersnum}{}
\renewcommand{\cftchapindent}{0cm}
\renewcommand{\cftchapnumwidth}{1.5em}
\renewcommand{\cftchappagefont}{\mdseries}
\renewcommand{\cftchapfont}{\mdseries}
\renewcommand{\cftchapleader}{\cftdotfill{\cftchapdotsep}}

% section entries
\renewcommand{\cftsecaftersnum}{}
\renewcommand{\cftsecindent}{1.5em}
\renewcommand{\cftsecnumwidth}{2.5em}
\renewcommand{\cftsecpagefont}{\mdseries}
\renewcommand{\cftsecfont}{\mdseries}
\renewcommand{\cftsecleader}{\cftdotfill{\cftsecdotsep}}

% subsection entries
\renewcommand{\cftsubsecaftersnum}{}
\renewcommand{\cftsubsecindent}{4em}
\renewcommand{\cftsubsecnumwidth}{3em}
\renewcommand{\cftsubsecpagefont}{\mdseries}
\renewcommand{\cftsubsecfont}{\mdseries}
\renewcommand{\cftsubsecleader}{\cftdotfill{\cftsubsecdotsep}}

% ============================================================

\newcommand{\reportAuthor}{Жуковін Олександр Віталійович}
\newcommand{\reportAuthorShort}{Олександр ЖУКОВІН}
\newcommand{\reportAuthorGroup}{ФІ-21}
\newcommand{\reportTitle}{Прогнозування інфляції України за допомогою методів машинного навчання}

\newcommand{\supervisorFio}{Чернятевич Антон Андрійович}
\newcommand{\supervisorFioShort}{Антон ЧЕРНЯТЕВИЧ}
\newcommand{\supervisorRegalia}{асистент кафедри ММАД}

\newcommand{\consultFio}{}
\newcommand{\consultRegalia}{}

\newcommand{\reviewerFio}{Прізвище Ім'я По-батькові}
\newcommand{\reviewerRegalia}{посада, степінь, звання}

\newcommand{\YearOfDefence}{2026}
\newcommand{\YearOfBeginning}{2025}

% ============================================================
\newcommand{\chapconclude}[1]{\section*{Висновки до розділу #1}\addcontentsline{toc}{section}{Висновки до розділу #1}}

\begin{document}

\pagestyle{plain}
\setfontsize{14}

% ============================================================
%  PAGE 1 — Title page
% ============================================================
\thispagestyle{empty}
\linespread{1.1}

\begin{center}
{\bfseries
НАЦІОНАЛЬНИЙ ТЕХНІЧНИЙ УНІВЕРСИТЕТ УКРАЇНИ \par
<<КИЇВСЬКИЙ ПОЛІТЕХНІЧНИЙ ІНСТИТУТ \par
імені Ігоря СІКОРСЬКОГО>>\par
Навчально-науковий фізико-технічний інститут\par
\medskip
Кафедра математичного моделювання та аналізу даних\par
}
\end{center}

\vspace{5mm}

\begin{tabularx}{\textwidth}{XX}
                       & <<До захисту допущено>> \\[06pt]
                       & в.о. завідувача кафедри \\[06pt]
                       & \rule{2.5cm}{0.25pt} Іван ТЕРЕЩЕНКО \\[06pt]
                       & <<\rule{0.5cm}{0.25pt}>> \rule{2.5cm}{0.25pt} \YearOfDefence~р.
\end{tabularx}

\begin{center}
\vspace{5mm}
{\bfseries\huge Дипломна робота} \par
{\bfseries на здобуття ступеня бакалавра} \par
за освітньо-професійною програмою \par
<<Математичні методи моделювання, розпізнавання образів та комп'ютерного зору>> \par
\end{center}

зі спеціальності: 113 Прикладна математика \par
на тему: \textbf{<<\reportTitle>>}

\vspace{5mm}

\begin{tabularx}{\textwidth}{>{\setlength\hsize{1.5\hsize}}X >{\setlength\hsize{0.5\hsize}}X}
Виконав:                                                         & \\
студент IV курсу, групи \reportAuthorGroup                       & \\
\reportAuthor                                                    & \rule{2.5cm}{0.25pt}   \\[12pt]
Керівник:                                                        & \\
\supervisorRegalia                                               & \\
\supervisorFio                                                   & \rule{2.5cm}{0.25pt}   \\[12pt]
Рецензент:                                                       & \\
\reviewerRegalia                                                 & \\
\reviewerFio                                                     & \rule{2.5cm}{0.25pt}
\end{tabularx}

\vspace{15mm}

\linespread{1.1}
\begin{tabularx}{\textwidth}{>{\setlength\hsize{1.25\hsize}}X >{\setlength\hsize{1.5\hsize}}X >{\setlength\hsize{0.25\hsize}}X}
& Засвідчую, що у цій дипломній роботі немає запозичень з праць інших
авторів без відповідних посилань.

& \\
& Студент \rule{2.5cm}{0.25pt}      &
\end{tabularx}

\vfill
\begin{center}
{Київ~-~\YearOfDefence}
\end{center}

\newpage
\thispagestyle{plain}

% ============================================================
%  PAGES 2–3 — Assignment / Specification
% ============================================================
\linespread{1.1}

\begin{center}
{\bfseries
НАЦІОНАЛЬНИЙ ТЕХНІЧНИЙ УНІВЕРСИТЕТ УКРАЇНИ \par
<<КИЇВСЬКИЙ ПОЛІТЕХНІЧНИЙ ІНСТИТУТ \par
імені Ігоря СІКОРСЬКОГО>>\par
Навчально-науковий фізико-технічний інститут\par
Кафедра математичного моделювання та аналізу даних}
\end{center}
\par

Рівень вищої освіти - перший (бакалаврський)

Спеціальність - 113~Прикладна математика,

ОПП <<Математичні методи моделювання, розпізнавання образів та комп'ютерного зору>>

\vspace{10mm}
\begin{tabularx}{\textwidth}{XX}
& ЗАТВЕРДЖУЮ                                      \\[06pt]
& Завідувач кафедри                               \\[06pt]
& \rule{2.5cm}{0.25pt} Іван ТЕРЕЩЕНКО             \\[06pt]
& <<\rule{0.5cm}{0.25pt}>> \rule{2.5cm}{0.25pt} \YearOfDefence~р.
\end{tabularx}

\vspace{5mm}
\begin{center}
{\bfseries ЗАВДАННЯ \par}
{\bfseries на дипломну роботу \par}
\end{center}

\frenchspacing
\doublespacing          % інтервал "1,5" між рядками, тепер навічно
\setfontsize{14}

Студент: \reportAuthor \par

1. Тема роботи: <<\emph{\reportTitle}>>,
науковий керівник дисертації: \supervisorRegalia~\supervisorFio, \par
затверджені наказом по університету \No~\rule{0.5cm}{0.25pt} від <<\rule{0.5cm}{0.25pt}>> \rule{2.5cm}{0.25pt} \YearOfDefence~р.

2. Термін подання студентом роботи: <<\rule{0.5cm}{0.25pt}>> \rule{2.5cm}{0.25pt} \YearOfDefence~р.

3. Вихідні дані до роботи: \emph{Макроекномічні показники України, Литви і Латвії, глобальні ринкові індикатори та індекси ринкових настроїв.}

4. Зміст роботи: \emph{Дослідження сучасних методів прогнозування інфляції. Побудова гібридної моделі на основі ARIMA та XGBoost методів й пояснювального ШІ. Валідація моделі та порівняння із традиційними методами.}

5. Перелік графічного матеріалу: \emph{Презентація доповіді}

6. Дата видачі завдання: 29 вересня \YearOfBeginning~р.

\newpage

\begin{center}
Календарний план
\end{center}

\renewcommand{\arraystretch}{1.5}
\begin{table}[h!]
\setfontsize{14pt}
\centering
    \begin{tabularx}{\textwidth}{|>{\centering\arraybackslash\setlength\hsize{0.25\hsize}}X|>{\setlength\hsize{2\hsize}}X|>{\centering\arraybackslash\setlength\hsize{1\hsize}}X|>{\centering\arraybackslash\setlength\hsize{0.75\hsize}}X|}
    \hline № \newline з/п & Назва етапів виконання дипломної роботи & Термін виконання & Примітка \\
    \hline
    1 &
    Узгодження теми роботи із науковим керівником. &
    вересень \YearOfBeginning~р. &
    Виконано \\
    \hline
    2 &
    Огляд опублікованих джерел за тематикою дослідження. &
    жовтень \YearOfBeginning~р. &
    Виконано \\
    \hline
    3 &
    Визначення необхідних показників та часових рамок даних для тренування та валідації моделі. &
    листопад \YearOfBeginning~р. &
    Виконано \\
    \hline
    4 &
    Пошук наборів даних макроекономічних показників України та світу. &
    грудень \YearOfDefence~р. &
    Виконано \\
    \hline
    5 &
    Збір необхідних даних показників та формування навчальної та тестової вибірок. &
    січень \YearOfDefence~р. &
    Виконано \\
    \hline
    6 &
    Аналіз існуючих математичних моделей прогнозування інфляції. &
    лютий \YearOfDefence~р. &
    Виконано \\
    \hline
    7 &
    Опис та побудова моделі з лінійною, нелінійною та пояснювальною складовими. &
    лютий--березень \YearOfDefence~р. &
    Виконано \\
    \hline
    8 &
    Проведення експерименту та оцінка результатів. &
    березень--квітень \YearOfDefence~р. &
    Виконано \\
    \hline
    9 &
    Оформлення результатів та підготовка презентації. &
    квітень--травень \YearOfDefence~р. &
    Виконано \\
    \hline
    \end{tabularx}
\end{table}

\renewcommand{\arraystretch}{1}
\begin{tabularx}{\textwidth}{>{\setlength\hsize{1.2\hsize}}X >{\setlength\hsize{0.5\hsize}}X >{\setlength\hsize{1.3\hsize}}X}
Студент  & \rule{2.5cm}{0.25pt}  & \reportAuthorShort \\[06pt]
Керівник & \rule{2.5cm}{0.25pt}  & \supervisorFioShort \\
\end{tabularx}

% ============================================================
%  РЕФЕРАТ
% ============================================================
\newpage
\thispagestyle{plain}
\setfontsize{14}

\begin{center}
    {\bfseries\Large РЕФЕРАТ}
\end{center}

\vspace{12pt}

Кваліфікаційна робота містить: \textbf{100}~стор., \textbf{100}~рисунків, \textbf{5}~таблиць, \textbf{27}~джерел.

У роботі досліджено сучасні методи прогнозування інфляції у центробанках держав та міжнародних фінансових інституціях. Для покращення результатів побудовано гібридну модель на основі методів ARIMA та XGBoost. Також для вирішення проблеми "чорного ящику" додано пояснювальну модель SHAP для інтерпретації результатів та визначення найбільш впливових чинників на інфляцію. Для тренувальної та валідаційної вибірок було використано макроекономічні показники України, Литви, Латвії, глобальні ринкові індикатори та індекси ринкових настроїв.

\vspace{6pt}
\MakeUppercase{інфляція, arima, xgboost, shap}

% ============================================================
%  ABSTRACT
% ============================================================
\newpage
\thispagestyle{plain}
\setfontsize{14}

\begin{center}
    {\bfseries\Large ABSTRACT}
\end{center}

\vspace{12pt}

The qualification work contains: \textbf{100}~pages, \textbf{100}~figures, \textbf{5}~tables, \textbf{27}~references.

This paper investigates modern methods of inflation forecasting used in national central banks and international financial institutions. To improve the results, a hybrid model based on ARIMA and XGBoost methods was developed. Additionally, to address the "black box" problem, the SHAP explanatory model was incorporated to interpret the results and identify the most influential factors driving inflation. Macroeconomic indicators of Ukraine, Lithuania, and Latvia, along with global market indicators and market sentiment indices, were utilized for the training and validation datasets.

\vspace{6pt}
\MakeUppercase{INFLATION, ARIMA, XGBOOST, SHAP}

\clearpage

% ============================================================
%  ЗМІСТ
% ============================================================
\tableofcontents

% ============================================================
%  ПЕРЕЛІК УМОВНИХ ПОЗНАЧЕНЬ, СИМВОЛІВ І СКОРОЧЕНЬ
% ============================================================
\newpage
\thispagestyle{plain}
\setfontsize{14}
\addcontentsline{toc}{chapter}{Перелік умовних позначень, символів і скорочень}

\begin{center}
    {\bfseries\Large\MakeUppercase{Перелік умовних позначень, символів і скорочень}}
\end{center}

\vspace{12pt}

{\setstretch{1.2}\setfontsize{14}
\noindent\begin{tabularx}{\textwidth}{@{}>{\bfseries}p{2.5cm}@{\hspace{0.5cm}}X@{}}
ARIMA    & Autoregressive Integrated Moving Average - авторегресійна інтегрована ковзна середня \\
CPI      & Consumer Price Index - індекс споживчих цін \\
LASSO    & Least Absolute Shrinkage and Selection Operator - метод мінімізації суми абсолютних значень коефіцієнтів \\
LSTM     & Long Short-Term Memory - довга короткострокова пам'ять \\
MAE      & Mean Absolute Error - середня абсолютна похибка \\
ML       & Machine Learning - машинне навчання \\
RF       & Random Forest - випадковий ліс \\
RMSE     & Root Mean Squared Error - корінь із середньоквадратичної похибки \\
SARIMA   & Seasonal Autoregressive Integrated Moving Average - сезонна авторегресійна інтегрована ковзна середня \\
SHAP     & SHapley Additive exPlanations - адитивні пояснення на основі значень Шеплі \\
VAR      & Vector Autoregression - векторна авторегресія \\
XAI      & Explainable Artificial Intelligence - пояснювальний штучний інтелект \\
XGBoost  & Extreme Gradient Boosting - метод екстремального градієнтного бустингу \\
ВВП      & валовий внутрішній продукт \\
НБУ      & Національний банк України \\
\end{tabularx}
\par}

% ============================================================
%  ВСТУП
% ============================================================
\newpage
\thispagestyle{plain}
\doublespacing
\setfontsize{14}
\addcontentsline{toc}{chapter}{Вступ}

\begin{center}
    {\bfseries\Large\MakeUppercase{Вступ}}
\end{center}

\vspace{12pt}

\textbf{Актуальність дослідження.}

У стабільний період на рівень інфляції діє безліч чинників, що зумовлює необхідність постійного коригування монетарної політики центральних банків задля її утримання в прийнятних межах. В умовах воєнного часу та в період післявоєнної відбудови ця проблема набуває особливої гостроти: показники інфляції можуть суттєво відхилятися від норми, що суттєво ускладнює економічне планування та знижує рівень життя населення. Тому ефективна оцінка інфляції є вкрай актуальним завданням, вирішення якого дозволить підвищити ефективність економічної політики держави.

\textbf{Метою дослідження} є прогнозування інфляції для України на основі внутрішніх та зовнішніх макроекономічних показників. Для досягнення мети необхідно дослідити наукові джерела на тематику роботи, зібрати дані для навчання та перевірки моделі, побудувати модель та натренувати її, провести експеримент та порівняти результат із традиційними методами.

\textbf{Об'єктом дослідження} є економіка України, а саме її макроекономічні показники та процеси, що визначають інфляційну динаміку в країні.

\textbf{Предметом дослідження} є кількісні економічні показники та процеси України, а також Латвії, Литви та світової економіки.

Дослідження передбачає збір та аналіз історичних економічних даних України, Латвії та Литви, встановлення кореляційних залежностей між зазначеними показниками, а також вивчення їхнього впливу на інфляційні процеси як в умовах економічної стабільності, так і в період криз. При розв'язанні поставлених завдань використовувались такі \textbf{методи дослідження}: методи математичної статистики та аналізу часових рядів, методи машинного навчання, методи ансамблевого моделювання, а також методи комп'ютерного моделювання та крос-валідації.

\textbf{Наукова новизна} ?

\textbf{Практичне значення} ?

% ============================================================
%  РОЗДІЛ 1
% ============================================================
\doublespacing
\setfontsize{14}

\chapter{Теоретичні основи інфляції та її прогнозування}

У цьому розділі, для розуміння економічних процесів що використовуються в роботі, розглядаються основи поняття інфляції, її причини, види та наслідки. Буде обґрунтована актуальність задачі прогнозування та описані складності що виникають при прогнозах. Також будуть наведені приклади сучасних класичних математичних моделей що використовуються для оцінки інфляції.

\section{Інфляція як економічне явище та її прояви в Україні}

У сучасній економіці інфляція та дефляція є фундаментальними поняттями, що описують протилежні за своїм характером зміни цін. Обидва процеси визначають купівельну спроможність грошей - кількість товарів або послуг, яку можна придбати за певну грошову одиницю. Інфляція супроводжується зростанням загального рівня цін і, відповідно, падінням купівельної спроможності грошей; дефляція, навпаки, означає зниження загального рівня цін, через що реальна вартість грошей зростає.

\subsection{Поняття інфляції та дефляції}


В понятті інфляції важливими є стабільність, тривалість і широта. Наприклад, зростання ціни на пальне через тимчасові перебої у постачанні - це звичайне ринкове коливання. Інфляція ж - це системний процес зростання загального рівня цін в економіці, що охоплює одночасно продовольчі товари, промислові вироби, послуги, комунальні тарифи, медикаменти, транспорт, освіту тощо. Саме системність та широта охоплення відрізняють інфляцію від звичайних цінових коливань окремих ринків~[1].

Дефляція, незважаючи на свою привабливість для споживачів, несе для економіки не менші ризики ніж інфляція. Коли споживачі очікують подальшого зниження цін, вони відкладають покупки, що скорочує попит. Виробники у відповідь скорочують виробництво, знижують заробітні плати або звільняють працівників. Падають прибутки компаній, зростають реальні борги (адже фактична сума боргу залишається незмінною, тоді як гроші дорожчають). Так формується <<дефляційна спіраль>>, що може призвести до глибокої рецесії~[5]. Класичним прикладом руйнівної дефляції є Велика депресія 1929--1933 років у США, коли загальний рівень цін впав майже на 25--30\%, що супроводжувалося масовим безробіттям та банківськими кризами.

Для кількісного вимірювання інфляції та дефляції в сучасній практиці використовують кілька основних показників. Найпоширенішим є індекс споживчих цін (CPI - Consumer Price Index), який відображає зміну вартості <<споживчого кошика>> - фіксованого набору товарів і послуг, що типово споживаються домогосподарствами~[1]. Формально темп інфляції за період~$t$ розраховується як відсоткова зміна індексу:

\begin{equation}
\pi_t = \frac{CPI_t - CPI_{t-1}}{CPI_{t-1}} \times 100\%
\end{equation}

де $\pi_t$ - темп інфляції за період~$t$ (у відсотках); $CPI_t$ - індекс споживчих цін у поточному періоді; $CPI_{t-1}$ - індекс споживчих цін у попередньому періоді. Найчастіше у звітах Національного банку України (НБУ) та інших центральних банків використовується річний показник інфляції - порівняння CPI поточного місяця з CPI того ж місяця попереднього року (year-over-year, або р/р).

Залежно від темпів зростання цін виділяють кілька типів інфляції. Помірна інфляція - це відносно повільне зростання цін, як правило до 10\% на рік або 1--2\% на місяць. Такий рівень традиційно вважається прийнятним для нормально функціонуючої ринкової економіки й навіть бажаним: гроші поступово знецінюються, тому їх вигідніше вкладати у виробничу діяльність, а не тримати в готівці. Більшість центральних банків розвинених країн встановлюють цільовий показник інфляції на рівні близько 2\% річних~[2].

Галопуюча інфляція передбачає значно вищі темпи зростання цін - від 10 до 100\% на рік. У цих умовах поведінка економічних суб'єктів кардинально змінюється: населення прагне якомога швидше позбутися грошей, купуючи товари тривалого використання або конвертуючи кошти в іноземну валюту; підприємства скорочують довгострокові інвестиції; контракти укладаються з прив'язкою до твердої валюти~[5].

Гіперінфляція - найбільш руйнівна форма знецінення грошей. У різних джерелах вона визначається по-різному: темпи понад 50\% на місяць, понад 100\% за три роки тощо. Спільним у всіх визначеннях є те, що гроші втрачають вартість настільки стрімко, що фактично перестають виконувати свої основні функції як засобу обігу і міри вартості~[2]. Сучасним прикладом є Венесуела, де з 2014 року триває економічна криза. Через різке падіння цін на нафту надходження в буджет значно впали що змусило владу збільшувати грошову масу шляхом друку коштів [6]. Динаміку цього процесу демонструє рис.~1.1.

\begin{figure}[ht]
\centering
\includegraphics[width=0.8\textwidth]{images/1.1.png}
\caption{Динаміка інфляції у Венесуелі в 2015--2017~рр. [6]}
\label{fig:weimar}
\end{figure}

Окремо варто розглянути поняття стагфляції - стану, за якого висока інфляція поєднується зі стагнацією економіки та зростанням безробіття. Стагфляція є особливо небезпечною, оскільки традиційні монетарні інструменти боротьби з інфляцією, наприклад підвищення відсоткових ставок, поглиблюють економічний спад, тоді як стимулюючі заходи для виходу з рецесії посилюють інфляцію~[5].

\subsection{Чинники, що впливають на інфляцію}

Інфляція є багатофакторним явищем, і розуміння її рушійних сил є ключовим як для економічного аналізу, так і для прогнозування. В економічній літературі прийнято виділяти кілька груп чинників: чинники попиту, чинники пропозиції, монетарні чинники, інфляційні очікування, а також зовнішні фактори~[5]. На практиці ці групи тісно взаємодіють між собою, і саме ця множинна взаємодія робить задачу прогнозування інфляції такою складною - лінійні моделі неспроможні адекватно її описати, що буде показано у підрозділі~1.2.

Інфляція попиту виникає, коли сукупний попит в економіці перевищує її виробничі можливості~[5]. Якщо рівень зайнятості та доходів населення зростає, споживачі мають більше грошей і прагнуть витратити їх на товари та послуги; коли виробнича система не здатна у короткостроковому періоді наростити пропозицію відповідними темпами, продавці підвищують ціни~[5].

Інфляція витрат виникає з боку пропозиції: зростання витрат виробництва змушує підприємства підвищувати ціни на кінцеву продукцію~[5]. Основними чинниками тут є зростання цін на енергоносії, підвищення вартості сировини та матеріалів, зростання заробітних плат, а також підвищення транспортних витрат~[5].

Монетарні чинники пов'язані з кількістю грошей в обігу та швидкістю їх обертання. Кількісна теорія грошей описується класичним рівнянням обміну Фішера:

\begin{equation}
M \times V = P \times Y
\end{equation}

де $M$ - грошова маса в обігу; $V$ - швидкість обертання грошей; $P$ - загальний рівень цін; $Y$ - реальний обсяг виробництва (реальний валовий внутрішній продукт, ВВП). За припущення про відносну стабільність~$V$ у короткостроковому періоді, темп зростання грошової маси, що перевищує темп зростання реального виробництва, у довгостроковому плані визначає темп інфляції~[5]. У крайньому випадку, коли уряд вдається до фінансування бюджетного дефіциту шляхом емісії, виникає ризик гіперінфляції. Саме тому в провідних країнах центральні банки є юридично незалежними від виконавчої влади установами~[2].

Інфляційні очікування - це суб'єктивні уявлення домогосподарств, підприємств, інвесторів тощо про майбутню інфляцію. Попри свій суб'єктивний характер, вони мають цілком реальний вплив на ціни: якщо підприємці очікують зростання цін на сировину, вони закладають це зростання у власні ціни заздалегідь; якщо працівники очікують інфляцію, вони вимагають підвищення заробітних плат, що знову ж таки збільшує витрати виробництва і підштовхує ціни~[5]. Так формується класичне <<самоздійснюване пророцтво>>. Управління інфляційними очікуваннями є одним із ключових завдань сучасних центральних банків.

Зовнішні чинники відіграють особливо значну роль у малих відкритих економіках, до яких належить й Україна. Зміна обмінного курсу є одним із найвпливовіших чинників: девальвація національної валюти автоматично підвищує ціни на імпортні товари, що прямо впливає на споживчий кошик. Цей механізм у літературі позначається як ефект перенесення і кількісно описується еластичністю:

\begin{equation}
\varepsilon = \frac{\Delta P/P}{\Delta E/E}
\end{equation}

де $\varepsilon$ - еластичність перенесення; $\Delta P/P$ - відносна зміна цін; $\Delta E/E$ - відносна зміна обмінного курсу. Для України, як і для більшості малих відкритих економік, $\varepsilon$ зазвичай знаходиться в діапазоні 0,2--0,4, що означає: 10-відсоткова девальвація гривні веде в середньому до 2--4-відсоткового зростання споживчих цін протягом наступних 6--12 місяців~[4].

\subsection{Інфляція в Україні}

Розглядаючи інфляційну історію України, важливо зосередитись на ключових подіях, що формують контекст для побудови прогностичних моделей.

Перший і найважчий період - гіперінфляція початку 1990-х. Лібералізація цін у 1992 році одразу виявила накопичені диспропорції: інфляція того року склала близько 2~000\%, а у 1993 році зросла до рекордних 10~256\%, що зробило Україну абсолютним рекордсменом серед країн, які не перебували у стані війни~[7]. Основним чинником була масштабна кредитна емісія НБУ для покриття бюджетного дефіциту, що досягав 12,2\% ВВП~[5].

\begin{figure}[ht]
\centering
\includegraphics[width=0.8\textwidth]{images/1.2.png}
\caption{Динаміка інфляції в Україні (тисячі \% р/р) з 1993 по 2000~рр. [4]}
\label{fig:ukraine_inflation_1}
\end{figure}

Другий період - час відносної стабільності 2000-х років (інфляція в межах 6--16\% на рік). На перший погляд, цей період міг би слугувати ідеальним матеріалом для побудови прогнозних моделей. Однак дослідження показують, що значною мірою стабільність забезпечувалася штучним утриманням обмінного курсу гривні на рівні близько 8~грн/дол. - кроком, що поступово накопичував зовнішні дисбаланси~[13]. Коли у 2014--2015 роках утримання курсу стало неможливим, корекція виявилась болісною: інфляція у 2015~р. досягла 48,7\%, а девальвація гривні склала близько 67\%~[11]. Цей епізод демонструє важливу властивість української інфляції - її чутливість до структурних змін у режимі обмінного курсу.

Третій період - час після початку повномасштабного вторгнення Росії у лютому 2022~р. Руйнування виробничих потужностей та логістичної інфраструктури різко скоротили пропозицію товарів, мільйони вимушених переселенців змінили географію попиту, масштабне фінансування воєнних витрат створило додатковий тиск на бюджет, а НБУ був змушений підвищити облікову ставку до 25\% - рекордного рівня~[13]. Інфляція у 2022~р. сягнула близько 26,6\%~[4].

\begin{figure}[ht]
\centering
\includegraphics[width=0.8\textwidth]{images/1.3.png}
\caption{Динаміка інфляції в Україні (\% р/р) з 2000 по 2024~рр. [4]}
\label{fig:ukraine_inflation_2}
\end{figure}

Узагальнюючи практичні уроки української інфляційної історії, можна сформулювати кілька ключових вимог до сучасної прогностичної моделі: (1)~здатність працювати в умовах нестаціонарності та структурних зрушень, (2)~урахування широкого спектру чинників, у тому числі курсу, монетарних агрегатів, цін на сировину, (3)~можливість швидкої адаптації до нових режимів економічної політики, (4)~інтерпретованість результатів - оскільки в умовах кризи довіра до моделі залежить від того, наскільки чітко її прогнози можна пояснити. Саме ці вимоги і обумовлюють перехід від традиційних до сучасних методів машинного навчання, доповнених інструментами пояснювального штучного інтелекту (ШІ).

\subsection{Антиінфляційна політика та таргетування інфляції в Україні}

Антиінфляційна політика держави і центрального банку є тим контекстом, у якому функціонує будь-яка прогностична модель. Розуміння цілей та інструментів цієї політики дозволяє правильно інтерпретувати прогнози та визначати їх практичну цінність. Для центрального банку точний прогноз інфляції - не самоціль, а основа для ухвалення рішень щодо монетарних інструментів. Якщо модель прогнозує, що інфляція протягом року перевищить цільовий орієнтир, регулятор має заздалегідь ужити заходів для її стримування.

Основним інструментом монетарної політики є ключова (облікова) ставка центрального банку. Підвищуючи її, регулятор робить запозичення дорожчими: банки отримують ресурси за вищою ціною і відповідно підвищують ставки за кредитами для бізнесу та населення~[2]. Це охолоджує попит на кредити, скорочує інвестиції та споживчі витрати, і зрештою знижує тиск на ціни. Зв'язок між реальною відсотковою ставкою та номінальною описується формулою Фішера:

\begin{equation}
r \approx i - \pi^e
\end{equation}

де $r$ - реальна відсоткова ставка; $i$ - номінальна відсоткова ставка (наприклад, облікова ставка центрального банку); $\pi^e$ - очікуваний рівень інфляції. Ця формула пояснює, чому центральні банки зазвичай встановлюють номінальну ставку вищою за поточну і прогнозовану інфляцію - щоб реальна ставка залишалась додатною і стимулювала заощадження у національній валюті~[9].

Ключовою інновацією у монетарній політиці стало запровадження режиму інфляційного таргетування. Суть підходу полягає у тому, що центральний банк публічно оголошує конкретний цільовий рівень інфляції і підпорядковує усю свою монетарну політику досягненню цієї мети~[2]. Інфляційне таргетування має кілька суттєвих переваг: публічний орієнтир дозволяє <<заякорювати>> інфляційні очікування, підвищує передбачуваність політики та захищає центральний банк від короткострокового політичного тиску.

В Україні режим інфляційного таргетування було запроваджено у 2015 році - у надзвичайно несприятливих умовах: країна перебувала у стані збройного конфлікту, економіка падала, банківська система потерпала від масового виведення депозитів, а гривня щойно пережила різку девальвацію~[8]. Тим не менш, реформа дала помітні результати: у 2017 році інфляція склала 13,7\%, у 2018-му - 9,8\%, у 2019-му - 4,1\%~[12], тобто вперше за довгий час опинилась нижче офіційної цілі НБУ.

Поточна ціль НБУ становить 5\% з допустимим інтервалом $\pm$1~п.п.~[3]. Згідно з Інфляційним звітом за січень 2026~р., монетарна політика спрямована на приведення інфляції до цієї цілі на горизонті, що не перевищує трьох років. Прогноз НБУ передбачає зниження інфляції до 7,5\% у 2026~р., до 6\% у 2027-му та досягнення цільових 5\% у 2028~р.~[3].

Однак дослідники зазначають, що механізми монетарної трансмісії в Україні залишаються відносно слабкими порівняно з розвиненими ринками~[10]. Це означає, що зміна облікової ставки не так швидко і повно відображається на кредитних ставках та сукупному попиті, як у країнах Європейського Союзу (ЄС). Однією з причин є низький рівень кредитування економіки приватним сектором. Слабкість механізмів трансмісії в українській економіці - додатковий аргумент на користь прогностичних моделей, що враховують широкий спектр немонетарних чинників: курсу, цін на енергоносії, бюджетних показників, зовнішніх шоків тощо. Саме такий мультифакторний підхід, що реалізується машинним навчанням, є основою даної роботи.

\section{Прогнозування та моделювання інфляції}

Розвиток кількісних методів прогнозування інфляції тривав упродовж усього XX~ст. і сформував широкий спектр інструментів, що нині позначаються як <<традиційні>> або <<класичні>>. До них відносять регресійні моделі, моделі авторегресії, векторні авторегресії, структурні макроекономічні моделі та моделі на основі кривої Філліпса. Кожен з цих підходів має власну теоретичну базу, переваги й обмеження, а також - що особливо важливо для контексту даної роботи - спільне принципове обмеження, яке і обумовлює перехід до сучасних методів.

Розвиток обчислювальних потужностей та зростання обсягів доступних даних призвели до використання методів машинного навчання (Machine Learning, ML) та глибокого навчання (Deep Learning, DL) для прогнозування інфляції. На відміну від традиційних підходів, що ґрунтуються на заздалегідь специфікованих лінійних залежностях, алгоритми ML здатні виявляти складні нелінійні патерни у великих масивах різнорідних даних без попереднього формулювання теоретичної структури моделі~[20].

\subsection{Традиційні підходи до прогнозування інфляції}

Одним з найбільш широко застосовуваних класів моделей є моделі авторегресії проінтегрованого ковзного середнього - ARIMA (Autoregressive Integrated Moving Average). Теоретичне підґрунтя цих моделей закладено в роботі Бокса і Дженкінса~[15], яка запропонувала систематичну методологію ідентифікації, оцінювання та верифікації моделей часових рядів. У загальному вигляді модель ARIMA описується рівнянням:

\begin{equation}
y_t = c + \varphi_1 y_{t-1} + \ldots + \varphi_p y_{t-p}
      + \varepsilon_t + \theta_1\varepsilon_{t-1} + \ldots + \theta_q\varepsilon_{t-q}
\end{equation}

де $y_t$ - поточне значення часового ряду; $c$ - константа; $\varphi_1,\ldots,\varphi_p$ - коефіцієнти авторегресійної складової порядку~$p$; $\theta_1,\ldots,\theta_q$ - коефіцієнти ковзного середнього порядку~$q$; $\varepsilon_t$ - випадкова похибка у момент~$t$. Параметр~$d$ позначає кількість взять різниць, необхідну для досягнення стаціонарності~[18].

Головною перевагою ARIMA є те, що для побудови моделі достатньо лише одного часового ряду - самого показника інфляції. Не потрібно збирати та підтримувати в актуальному стані великий масив допоміжних даних. Дослідження ірландської, нігерійської та суданської інфляції підтвердили, що ARIMA-моделі дають задовільні результати на коротких горизонтах прогнозування - до трьох-шести місяців - у стабільних економічних умовах~[14,~16]. Однак ARIMA повністю ігнорує будь-які зовнішні чинники - рух обмінного курсу, зміни у грошовій масі, цінові шоки на сировинних ринках. Якщо відбуваються структурні зрушення, прогнози різко погіршуються~[14].

Розширенням ARIMA є модель SARIMA (Seasonal ARIMA), що додатково включає сезонні компоненти. Для інфляції, яка часто демонструє стійкі сезонні патерни (наприклад, через урожаї або тарифні перегляди), SARIMA є більш точною. На практиці SARIMA є одним із найбільш стійких бенчмарків при прогнозуванні споживчих цін і широко застосовується у дослідженнях центральних банків~[14].

Принципово іншим класом є векторні авторегресійні моделі - VAR (Vector Autoregression). На відміну від ARIMA, що є моделлю одного ряду, VAR одночасно моделює систему взаємозалежних змінних. У системі VAR кожна змінна описується як лінійна функція власних минулих значень та минулих значень усіх інших змінних моделі. Для прогнозування інфляції типовими змінними VAR є: інфляція, ВВП, відсоткова ставка, обмінний курс, грошова маса, ціни на нафту~[19].

VAR-модель здатна враховувати взаємозворотні причинно-наслідкові зв'язки між змінними і є потужним інструментом для сценарного аналізу. Водночас класичні VAR-моделі страждають на <<прокляття розмірності>>: з ростом числа змінних та лагів кількість параметрів, що підлягають оцінюванню, зростає квадратично~[18]. Відповіддю на цю проблему стала байєсівська векторна авторегресія (BVAR - Bayesian Vector Autoregression). BVAR-моделі широко застосовуються центральними банками, у тому числі Федеральною резервною системою США (ФРС) та Європейським центральним банком (ЄЦБ)~[14].

Окремий важливий клас становлять моделі на основі кривої Філліпса, яка відображає залежність між рівнем безробіття та рівнем інфляції. Пізніші розробки - праці Фрідмана і Фелпса у 1960-х - ввели в модель роль інфляційних очікувань. Однак починаючи з 1990-х років прогнози суттєво погіршали, адже зв'язок між безробіттям та інфляцією значно ослаб (так зване «сплющення» кривої Філліпса)~[14].

Найбільш структурно завершеним класом є динамічні стохастичні моделі загальної рівноваги - DSGE (Dynamic Stochastic General Equilibrium). DSGE-моделі ґрунтуються на мікроекономічних засадах і є <<серцем>> прогнозних систем більшості великих центральних банків~[18]. Вони мають сильну теоретичну інтерпретованість, проте вимагають значної кількості теоретичних припущень і нерідко поступаються BVAR у точності~[14].

Спільне принципове обмеження більшості традиційних методів - лінійність припущень. Вони погано вловлюють нелінійні взаємозалежності, що є характерними для реальних інфляційних процесів, особливо під час криз і структурних зрушень~[20]. Для України, де структурні зрушення є радше правилом, ніж винятком, це обмеження є особливо критичним. Саме воно і зумовлює необхідність переходу до сучасних методів машинного навчання.

\subsection{Сучасні підходи: машинне навчання у прогнозуванні інфляції}

Серед алгоритмів ML особливе місце посідають ансамблеві методи на основі дерев рішень. Метод випадкового лісу (Random Forest, RF) будує велику кількість дерев рішень на різних підвибірках даних і різних підмножинах ознак, а кінцевий прогноз формується шляхом усереднення. Така агрегація дозволяє суттєво знизити дисперсію прогнозу і уникнути перенавчання. Дослідження Медейруса~[21] показало, що Random Forest перевершує традиційні бенчмарки при прогнозуванні інфляції в США на основі понад 100 макроекономічних змінних, особливо у нестабільні економічні періоди.

Градієнтний бустинг (Gradient Boosting, GB) та його найвідоміша реалізація - XGBoost (Extreme Gradient Boosting) - є ще одним потужним ансамблевим методом. На відміну від RF, що будує дерева паралельно, бустинг навчає їх послідовно: кожне наступне дерево зосереджується на прикладах, де попередня ансамблева модель помилилась найбільше~[22]. Дослідницька робота Міжнародного валютного фонду (МВФ) 2024~р. <<Mending the Crystal Ball>> демонструє, що ансамблеві ML-методи здатні суттєво покращити прогнозування базової інфляції на горизонтах від 1 до 6~місяців. При цьому якість прогнозу особливо зростає у ті моменти, коли стандартні економічні залежності порушуються через нестандартні шоки - саме там, де традиційні методи найбільш вразливі~[22].

Регуляризовані лінійні моделі - LASSO (Least Absolute Shrinkage and Selection Operator), Ridge-регресія та їхнє поєднання Elastic Net - займають проміжне місце між традиційними лінійними методами та складнішими алгоритмами ML. LASSO мінімізує функцію втрат із додаванням штрафу за $L_1$-норму вектора коефіцієнтів:

\begin{equation}
L(\beta) = \sum_{i=1}^{n}(y_i - \hat{y}_i)^2 + \lambda\sum_{j=1}^{p}|\beta_j|
\end{equation}

де $\beta$ - вектор коефіцієнтів; $y_i$,~$\hat{y}_i$ - спостережене та прогнозоване значення; $\lambda$ - параметр регуляризації, що контролює силу штрафу. Завдяки $L_1$-штрафу LASSO автоматично обнуляє менш важливі предиктори і тим самим виконує відбір змінних~[22]. При великій кількості потенційних предикторів LASSO є ефективним інструментом приведення моделі до найбільш інформативної підмножини. Деякі дослідження показують, що LASSO-регресія нерідко конкурентоспроможна з глибокими нейромережевими моделями, особливо на коротших горизонтах прогнозування~[14].

Нейронні мережі і глибоке навчання відкривають широкі можливості для моделювання нелінійних залежностей. Особливу роль у роботі з часовими рядами відіграє архітектура рекурентних нейронних мереж, зокрема мереж з довготривалою і короткотривалою пам'яттю - LSTM (Long Short-Term Memory). Завдяки спеціальній архітектурі комірок пам'яті з вентилями LSTM здатна зберігати і передавати інформацію через великі часові проміжки~[14]. Однак практичні результати застосування LSTM до прогнозування інфляції є неоднозначними. Ряд досліджень фіксує, що LSTM поступається простішим моделям - SARIMA, Random Forest або навіть LASSO - особливо на коротких горизонтах~[14].

Останніми роками увагу привертають архітектури на основі механізму уваги і трансформери. На відміну від LSTM, що обробляє послідовність покроково, трансформер розглядає всю послідовність одразу і явно зважує взаємозалежності між різними часовими точками~[16]. Дослідження прогнозування місячної інфляції США показало, що трансформер демонструє найнижчі показники помилок RMSE і MAE серед усіх порівнюваних моделей.

Стандартними метриками для оцінки якості прогнозних моделей є середньоквадратична помилка (Root Mean Squared Error, RMSE) та середня абсолютна помилка (Mean Absolute Error, MAE):

\begin{equation}
\begin{aligned}
RMSE &= \sqrt{\frac{1}{n}\sum_{i=1}^{n}(y_i - \hat{y}_i)^2} \\[6pt]
MAE  &= \frac{1}{n}\sum_{i=1}^{n}|y_i - \hat{y}_i|
\end{aligned}
\end{equation}

де $y_i$ - фактичне значення; $\hat{y}_i$ - прогнозоване значення; $n$ - кількість спостережень у позавибірковому періоді. На практиці зазвичай використовуються обидва показники одночасно для більш повної оцінки.

Сучасні підходи становлять широкий спектр методів - від регуляризованих лінійних моделей через ансамблеві деревні методи до глибокого навчання та трансформерів. Жоден з цих методів не є безумовно найкращим у всіх ситуаціях: вибір залежить від горизонту прогнозу, обсягу і якості даних, наявності структурних зрушень. Однак, як показують дослідження, ансамблеві деревні методи (особливо XGBoost та Random Forest) є найбільш збалансованим вибором для задач макроекономічного прогнозування з обмеженим обсягом даних - поєднуючи високу точність, відносну стабільність та можливість інтерпретації результатів.

\subsection{Пояснювальний штучний інтелект та метод SHAP}

Висока точність прогнозів сама по собі не вирішує всіх проблем - без можливості пояснити, чому модель видає той чи інший результат, її практична цінність для прийняття рішень дуже обмежена. Особливо це стосується макроекономічного прогнозування, адже центральний банк, що ухвалює рішення про підвищення облікової ставки, має обґрунтувати своє рішення. Саме тому пояснювальний штучний інтелект (Explainable AI, XAI) є невід'ємним компонентом сучасного підходу до моделювання інфляції.

Серед методів XAI на сьогодні найпотужнішим і найбільш теоретично обґрунтованим є SHAP (SHapley Additive exPlanations), запропонований Скоттом Лундбергом і Су-Ін Лі у 2017~р.~[23]. SHAP базується на концепції значень Шеплі - підході з теорії кооперативних ігор. Якщо розглядати ознаки моделі як <<гравців>>, що спільно роблять <<прогноз-виграш>>, то значення Шеплі вказує, яким повинен бути справедливий розподіл цього виграшу між гравцями.

Формально, SHAP-значення ознаки~$i$ для конкретного спостереження визначається як:

\begin{equation}
\varphi_i = \sum_{S\subseteq N\setminus\{i\}}
\frac{|S|!\,(|N|-|S|-1)!}{|N|!}
\bigl[f(S\cup\{i\})-f(S)\bigr]
\end{equation}

де $\varphi_i$ - SHAP-значення ознаки~$i$; $N$ - повна множина ознак моделі; $S$ - підмножина ознак, що не містить~$i$; $|S|$ та~$|N|$ - розміри відповідних множин; $f(S)$ - прогноз моделі при використанні лише ознак з~$S$~[24]. Інтуїтивно ця формула описує середній маржинальний внесок ознаки~$i$, обчислений по всіх можливих підмножинах інших ознак.

SHAP надає два рівні пояснень - локальний та глобальний. На локальному рівні для кожного окремого прогнозу можна обчислити, який внесок зробила кожна ознака. На глобальному рівні SHAP-значення можна агрегувати по всіх спостереженнях, щоб зрозуміти, які ознаки в цілому є найбільш важливими і як вони впливають на прогнози~[26].

Технічно для деревних ансамблевих моделей (Random Forest, XGBoost, LightGBM) існує спеціалізована ефективна реалізація - TreeSHAP, що дозволяє обчислювати точні SHAP-значення з обчислювальною складністю $O(TLD^2)$ замість $O(TL \cdot 2^M)$ для загального випадку, де $T$ - кількість дерев, $L$ - максимальна кількість листків, $D$ - глибина, $M$ - кількість ознак~[24]. Ця ефективність робить SHAP практично застосовним навіть для великих ансамблевих моделей з десятками ознак.

У контексті прогнозування інфляції SHAP має додаткові важливі переваги. По-перше, він дозволяє відстежувати у часі, як змінюється відносна важливість різних чинників: у різні економічні режими ключову роль можуть відігравати різні змінні. По-друге, SHAP-аналіз можна використовувати як інструмент сценарного моделювання, відповідаючи на питання <<що станеться з прогнозом, якщо облікова ставка зросте на 1~п.п.>>~[25]. Це робить його не просто методом інтерпретації, а повноцінним аналітичним інструментом для підтримки прийняття рішень.

Поєднання сучасних точних методів машинного навчання з прозорим поясненням робить прогнозування якісно новим у порівнянні як з традиційними економетричними моделями, так і з важкоінтерпретованими підходами ML.

\subsection{Дані: проблема обмеженості та залучення країн-аналогів}

Якість будь-якої прогностичної моделі, особливо побудованої на методах машинного навчання, визначається не лише вибором алгоритму, а й якістю та обсягом доступних даних. Для української макроекономіки питання даних є одним з найбільш проблемних.

Перша і найбільша проблема - обмежений часовий горизонт однорідних даних. Як було показано у підрозділі~1.1.3, інфляційний ряд України до 1996 року є аномальним через гіперінфляцію і фактично не може бути використаний для побудови моделей, що мають прогнозувати сучасні значення. Період 1996--2014 рр. є відносно стабільним, але характеризується режимом фіксованого обмінного курсу, який після 2014 року було замінено інфляційним таргетуванням. Таким чином, реально придатний для побудови сучасних ML-моделей період охоплює близько 10--11 років, або приблизно 120--130 місячних спостережень~[12]. Для класичних статистичних моделей такий обсяг є достатнім, але для алгоритмів машинного навчання він є на межі мінімально допустимого.

Друга проблема - нестаціонарність та структурні зрушення у самих даних. Серед цих 120--130 спостережень частина припадає на період пандемії COVID-19 (2020--2021~рр.), частина - на повномасштабну війну (з 2022~р.), що супроводжується безпрецедентними шоками всіх типів. Кожен такий епізод фактично створює <<новий режим>>, що знижує ефективність навчання моделі на історичних даних попередніх режимів.

Третя проблема - обмежений набір предикторів з достатньою глибиною. Для багатьох важливих макроекономічних показників в Україні часові ряди починаються лише з 2000-х або навіть 2010-х рр., а ряд інших публікуються з затримками або переглядаються. Деякі важливі джерела (наприклад, дані про настрої споживачів, високочастотні індикатори ділової активності) у повному обсязі доступні лише з кінця 2010-х.

В такому випадку для подолання обмеженості даних є логічним збільшення вибірки за рахунок включення схожих спостережень. У контексті макроекономіки це означає залучення даних інших країн, що мають подібні структурні характеристики. У даній роботі такими країнами-аналогами обрано Латвію та Литву. Цей вибір обумовлений кількома взаємопов'язаними причинами.

Латвія, Литва та Україна пройшли спільний шлях трансформації від радянської командної економіки до ринкової. Усі три країни на початку 1990-х пережили періоди екстремально високої інфляції: Латвія - 958,6\% у 1992~р., Литва - близько 1~100\% у 1992-му, Україна - 10~256\% у 1993-му~[24,~25]. Усі три відновили самостійні центральні банки, провели грошові реформи (Литва ввела літас у 1993~р., Латвія - лат у 1993-му, Україна - гривню у 1996-му). Це означає принципову подібність структурних викликів у монетарній сфері.

Латвія та Литва, як і Україна, є малими відкритими економіками з високою чутливістю до зовнішніх шоків - насамперед курсових та сировинних. У всіх трьох країнах ефект перенесення обмінного курсу відіграє ключову роль у формуванні інфляційної динаміки~[27].

У 2022~р., коли Україна зіткнулася з повномасштабним російським вторгненням, інфляція в Латвії досягла 17,2\%, у Литві - 18,9\%~[27]. Хоча природа шоків була різною (для України - пряма війна, для країн Балтії - енергетичні наслідки відмови від російських енергоносіїв та геополітичні ризики), макроекономічна динаміка в усіх трьох країнах продемонструвала якісно подібні патерни: різке зростання цін на енергоносії, девальваційний (для України) або інфляційний (для країн єврозони) тиск, жорсткість монетарної політики у відповідь.

Латвія і Литва мають значно довші та якісніші ряди макроекономічних даних. Їхнє членство в ЄС (з 2004~р.) та єврозоні (Латвія з 2014~р., Литва з 2015~р.) забезпечило стандартизацію статистики за методологією Eurostat, тривалу історію інфляційного таргетування та повне розкриття даних. Це дозволяє суттєво розширити навчальну вибірку, додавши до приблизно 130 українських спостережень ще близько 250--300 спостережень з кожної з двох балтійських країн.

Технічно об'єднання даних трьох країн виконується у формі панельного ML-датасету, де кожне спостереження характеризується не лише економічними показниками, але й ідентифікатором країни як категоріальною ознакою. Це дозволяє моделі фіксувати спільні закономірності, що діють у всіх трьох економіках, та водночас враховувати специфіку кожної.

\chapconclude{1}

У цьому розділі було обґрунтовано актуальність проблеми інфляції у всі часи та складність прогнозування цього економічного процесу через його мультифакторність та чутливість до великої кількості чинників. Були розглянуті класичні підходи прогнозування, зокрема ARIMA моделі, та сучасні ML рішення, зокрема XGBoost. Також було обґрунтовано важливість впровадження XAI для практичного використання сучасних моделей з використанням методів машинного навчання для розуміння ключових факторів що мали вплив на той чи інший прогноз та винесення відповідного рішення для стримування інфляції.

% ============================================================
%  Список використаних джерел (ВИПРАВЛЕНА ВЕРСІЯ)
% ============================================================
\newpage
\addcontentsline{toc}{chapter}{Список використаних джерел}
\begin{center}
    {\bfseries\Large\MakeUppercase{Список використаних джерел}}
\end{center}

\vspace{12pt}

\begin{enumerate}[label={[\arabic*]}, leftmargin=0pt, itemindent=0pt, labelindent=0pt]

\item Measuring inflation and consumer prices [Електронний ресурс] / European Central Bank. --- URL: \url{https://www.ecb.europa.eu/stats/macroeconomic_and_sectoral/hicp/html/index.en.html}.

\item Lagarde C. Monetary policy in a high inflation environment: commitment and clarity [Електронний ресурс] : speech at the 5th annual conference organised by the European Fiscal Board, Frankfurt, 4 November 2022 / C.~Lagarde // European Central Bank. --- URL: \url{https://www.ecb.europa.eu/press/key/date/2022/html/ecb.sp221104_1~8be9a4f4c1.en.html}.

\item Національний банк України. Інфляційний звіт. Січень 2026 року [Електронний ресурс]. --- К.~: НБУ, 2026. --- URL: \url{https://bank.gov.ua/admin_uploads/article/IR_2026-Q1.pdf}.

\item VoxUkraine. Inflation in Ukraine: Past, Present and Future [Електронний ресурс]. --- URL: \url{https://voxukraine.org/en/inflation-in-ukraine-past-present-and-future}.

\item Кириленко В.~В. Макроекономіка: курс лекцій [Електронний ресурс] / за ред. В.~В.~Кириленка. --- Тернопіль~: Економічна думка, 2008. --- 252~с. --- URL: \url{https://dspace.wunu.edu.ua/handle/316497/28895}.

\item Brunnermeier M.~K. The Debt-Inflation Channel of the German Hyperinflation [Електронний ресурс] / M.~K.~Brunnermeier, S.~Correia, S.~Luck, E.~Verner, T.~Zimmermann // MPRA Paper. --- 2023. --- No.~119195. --- URL: \url{https://mpra.ub.uni-muenchen.de/119195/1/MPRA_paper_119195.pdf}.

\item Banaian K. The Ukrainian Economy since Independence [Електронний ресурс] / K.~Banaian. --- Cheltenham~: Edward Elgar Publishing, 1999. --- 144~p. --- (Studies of Communism in Transition). --- URL: \url{https://www.e-elgar.com/shop/usd/the-ukrainian-economy-since-independence-9781858989907.html}.

\item Kravchuk R.~S. Budget Deficits, Hyperinflation and Stabilization in Ukraine, 1991--96 / R.~S.~Kravchuk // Public Budgeting~\& Finance. --- 1998. --- Vol.~18(4). --- P.~45--70. --- DOI:~10.1046/j.0275-1100.1998.01149.x. --- URL: \url{https://doi.org/10.1046/j.0275-1100.1998.01149.x}.

\item Zholud O. The Effectiveness of the Monetary Transmission Mechanism in Ukraine since the Transition to Inflation Targeting / O.~Zholud, V.~Lepushynskyi, S.~Nikolaychuk // Visnyk of the National Bank of Ukraine. --- 2019. --- No.~247. --- P.~19--37. --- DOI:~10.26531/vnbu2019.247.02. --- URL: \url{https://doi.org/10.26531/vnbu2019.247.02}.

\item Grui A. Wartime Interest Rate Pass-through in Ukraine: The Role of Prudential Policy [Електронний ресурс] / A.~Grui // IES Working Paper. --- 2024. --- No.~33/2024. --- Charles University in Prague. --- URL: \url{https://ies.fsv.cuni.cz/sites/default/files/uploads/files/wp_2024_33_grui_2.pdf}.

\item Faryna O. Nonlinear Exchange Rate Pass-Through to Domestic Prices in Ukraine / O.~Faryna // Visnyk of the National Bank of Ukraine. --- 2016. --- No.~236. --- P.~30--42. --- DOI:~10.26531/vnbu2016.236.030. --- URL: \url{https://doi.org/10.26531/vnbu2016.236.030}.

\item A Decade of Inflation Targeting in Ukraine: Lessons for a New Era of Uncertainty [Електронний ресурс] // VoxUkraine. --- 2025. --- URL: \url{https://voxukraine.org/en/a-decade-of-inflation-targeting-in-ukraine-lessons-for-a-new-era-of-uncertainty}.

\item Rakic D. Two years of war: The state of the Ukrainian economy in ten charts [Електронний ресурс] : Briefing PE~747.858 / D.~Rakic // Economic Governance and EMU Scrutiny Unit, European Parliament. --- 2024. --- URL: \url{https://www.europarl.europa.eu/RegData/etudes/BRIE/2024/747858/IPOL_BRI(2024)747858_EN.pdf}.

\item Paranhos L. Inflation Forecasting: LSTM Networks vs.\ Traditional Models for Accurate Predictions / L.~Paranhos // Journal of Risk and Financial Management. --- 2025. --- Vol.~18(7). --- Art.~365. --- DOI:~10.3390/jrfm18070365. --- URL: \url{https://doi.org/10.3390/jrfm18070365}.

\item Meyler A. Forecasting Irish Inflation Using ARIMA Models [Електронний ресурс] / A.~Meyler, G.~Kenny, T.~Quinn // Central Bank of Ireland Technical Paper. --- 1998. --- No.~3/RT/98. --- URL: \url{https://mpra.ub.uni-muenchen.de/11359/}.

\item Nortey E.~N.~N. AI meets economics: Can deep learning surpass machine learning and traditional statistical models in inflation time series forecasting? / E.~N.~N.~Nortey [et al.] // Data Science in Finance and Economics. --- 2025. --- Vol.~5(2). --- P.~136--155. --- DOI:~10.3934/DSFE.2025007. --- URL: \url{https://doi.org/10.3934/DSFE.2025007}.

\item Modelling and Forecasting Inflation Rates in Kenya Using ARIMA Models [Електронний ресурс] // European Journal of Statistics and Probability. --- 2023. --- Vol.~11(1). --- P.~54--68. --- URL: \url{https://eajournals.org/ejsp/vol11-issue1-2023/modelling-and-forecasting-inflation-rates-in-kenya-using-arima-model/}.

\item Mwangangi A.~K. An Innovative Approach to Short-Term Inflation Forecasting in Rwanda [Електронний ресурс] / A.~K.~Mwangangi // BNR Economic Review. --- Vol.~21(2). --- URL: \url{https://www.ajol.info/index.php/BNRER/article/view/290174}.

\item Short-Term Inflation Forecasting in Sierra Leone: A Comparison of VAR, ARIMAX, and ARIMA Models // International Journal of Scientific Research and Management. --- 2025. --- DOI:~10.18535/ijsrm/v13i05.em04. --- URL: \url{https://doi.org/10.18535/ijsrm/v13i05.em04}.

\item Ismail L.~B. Modeling inflation with machine learning: a cross-horizon systematic review / L.~B.~Ismail [et al.] // International Journal of Data Science and Analytics. --- 2025. --- DOI:~10.1007/s41060-025-00905-w. --- URL: \url{https://doi.org/10.1007/s41060-025-00905-w}.

\item Naghi A.~A. The benefits of forecasting inflation with machine learning: New evidence / A.~A.~Naghi, E.~O'Neill, M.~Medeiros // Journal of Applied Econometrics. --- 2024. --- DOI:~10.1002/jae.3088. --- URL: \url{https://doi.org/10.1002/jae.3088}.

\item IMF. Mending the Crystal Ball: Enhanced Inflation Forecasts with Machine Learning [Електронний ресурс] // IMF Working Paper. --- 2024. --- No.~WP/24/206. --- URL: \url{https://www.imf.org/-/media/Files/Publications/WP/2024/English/wpiea2024206-print-pdf.ashx}.

\item Lundberg S.~M. A Unified Approach to Interpreting Model Predictions [Електронний ресурс] / S.~M.~Lundberg, S.-I.~Lee // Advances in Neural Information Processing Systems (NeurIPS). --- 2017. --- Vol.~30. --- URL: \url{https://arxiv.org/abs/1705.07874}.

\item SHAP analysis: Explaining supervised machine learning model predictions [Електронний ресурс] // CPT: Pharmacometrics~\& Systems Pharmacology. --- 2024. --- URL: \url{https://pmc.ncbi.nlm.nih.gov/articles/PMC11513550/}.

\item Kn\"obl A. Stabilization in the Baltic Countries: Early Experience [Електронний ресурс] / A.~Kn\"obl, R.~Haas // Road Maps of the Transition: The Baltics, the Czech Republic, Hungary, and Russia. IMF Occasional Paper. --- 1996. --- No.~127. --- Washington, DC~: International Monetary Fund. --- URL: \url{https://www.elibrary.imf.org/display/book/9781557755193/ch01.xml}.

\item Bit\={a}ns M. Economic and Monetary Developments in Latvia (1945--1991) [Електронний ресурс] / M.~Bit\={a}ns, V.~Purvi\c{n}\v{s} // Bank of Latvia. --- URL: \url{https://www.bank.lv/images/stories/pielikumi/publikacijas/citaspublikacijas/Bitans_Purvins_EN.pdf}.

\item Economic Evolution in Euro-Adopting States vs.\ Future Adopters: A Comparative Analysis [Електронний ресурс] // Economies (MDPI). --- 2025. --- Vol.~13(8). --- Art.~239. --- URL: \url{https://www.mdpi.com/2227-7099/13/8/239}.

\end{enumerate}

\end{document}
